%------------------------------------------------------------------------------%
\section{Função composta}

%------------------------------------------------------------------------------%
\begin{definicao}{Função composta}{def:funcao_composta}
  Dadas as funções $f$ e $ g$ cujos domínios são $A$ e $B$, respectivamente,
  definimos a \emph{função composta}\index{Função!composta} $f \circ g$ por
  \[
    (f \circ g)(x)= f(g(x))
  \]
  O domínio de $f \circ g$ é o conjunto dos valores de $x \in B$
  tais que $g(x) \in A$
\end{definicao}
%------------------------------------------------------------------------------%

%------------------------------------------------------------------------------%
\begin{example}
  Se $f(x)= x^2$ e $g(x)=x-3$, encontre as compostas
  $f \circ g$ e $g \circ f$.

  \solution
  Por definição
  $f \circ g= f(g(x))=f(x-3)=(x-3)^2$ e
  $g \circ f= g(f(x))=g(x^2)= x^2-3$.
\end{example}
%------------------------------------------------------------------------------%

%------------------------------------------------------------------------------%
\begin{example}
  Se
  $f(x) = \frac{x}{x+1}$,
  $g(x) = x^{10}$ e
  $h(x) = x+3$,
  encontre  $f \circ g \circ h$.

  \solution
  \[
    f \circ g \circ h
    = f(g(h(x)))
    = f(g(x+3))
    = f((x+3)^{10})
    = \frac{(x+3)^{10}}{(x+3)^{10}+1}.
  \]
\end{example}
%------------------------------------------------------------------------------%

%------------------------------------------------------------------------------%
\begin{observation}
  Devemos está atentos ao calcularmos o domínio da função composta. Segundo a
  definição~\ref{def:funcao_composta}
  para que sejamos capazes de calcular $f(g(x))$ é preciso que:
  \begin{enumerate}
    \item $x$ pertença ao domínio de $g$;
    \item $g(x)$ (o valor de $g$ em $x$) pertença ao domínio de $f$;
  \end{enumerate}
\end{observation}
%------------------------------------------------------------------------------%

Para deixar mais clara a definição do domínio da função composta, veremos
como obtê-lo de forma prática no exemplo a seguir.

%------------------------------------------------------------------------------%
\begin{example}
  Sejam dadas as funções
  \[
    f(x) = \frac{1}{x^2-4} \quad\text{ e }\quad
    g(x) = \frac{1}{x}.
  \]
  Obtenha $f(g(x))$ e determine o domínio dessa composta.

  \solution
  \[
    f(g(x))
    = f\left(\frac{1}{x}\right)
    = \frac{1}{\dfrac{1}{x^2}-4}
    = \frac{1}{\dfrac{1-4x^2}{x^2}}
    = \frac{x^2}{1-4x^2}.
  \]
  Pela Observação,
  para que $f(g(x))$ esteja bem definida temos que ter:
  \begin{enumerate}
    \item $x$ pertença ao domínio de $g$, isto é, $x \neq 0$.
    \item $g(x)$ pertença ao domínio de $f$, isto é
    \[
      1-4x^2 \neq 0
      \quad\Rightarrow\quad
      x \neq -\frac{1}{2}
      \quad\text{ e }\quad
      x \neq \frac{1}{2}.
    \]
  \end{enumerate}
  Logo
  \[
    D(f(g(x)))
    = \left\{
        x \in \R \st
        x \neq 0,\;
        x \neq -\frac{1}{2},\;
        x \neq  \frac{1}{2}
    \right\}
  \]
\end{example}
%------------------------------------------------------------------------------%

%------------------------------------------------------------------------------%